\section*{Capítulo \ref{ch:b3:surfaces-catalysis} --- Superficies, adsorción y catálisis heterogénea}

\begin{solution}{exo:b3:surfaces-catalysis:1}
$K = \theta/((1 - \theta)p) = 0.40/(0.60 \times 2.0) = \qty{0.33}{kPa^{-1}}$; a \qty{10}{kPa}, $\theta = 3.33/4.33 = 0.77$.
\end{solution}

\begin{solution}{exo:b3:surfaces-catalysis:2}
$-20$: \omterm{def:b3:surfaces-catalysis:physi-chemi}{fisisorción} (del orden de una entalpía de condensación); $-150$: \omterm{def:b3:surfaces-catalysis:physi-chemi}{quimisorción}, la única
que puede ser disociativa.
\end{solution}

\begin{solution}{exo:b3:surfaces-catalysis:3}
(100): un sitio superior, dos puente (cuatro enlaces compartidos por dos átomos cada uno), un hueco de
cuatro (cuatro huecos compartidos por cuatro átomos). (111): un sitio superior, tres puente, dos
huecos de tres (seis triángulos en torno a cada átomo, cada uno compartido por tres).
\end{solution}

\begin{solution}{exo:b3:surfaces-catalysis:4}
$\num{3.0e-6}/\num{1.5e-5} = \qty{0.20}{s^{-1}}$.
\end{solution}

\begin{solution}{exo:b3:surfaces-catalysis:5}
$\theta/(1 - \theta) = (Kp)^{1/2}$: $0.111$ a \qty{1.0}{kPa}, el doble a \qty{4.0}{kPa}: $0.222$, luego $\theta = 0.18$.
\end{solution}

\begin{solution}{exo:b3:surfaces-catalysis:6}
Denominador: $1 + 10 + 2.5 = 13.5$: $\theta_{\ce{CO}} = 0.74$, $\theta_{\ce{O2}} = 0.19$, sitios libres 0.07. La
velocidad de Langmuir--Hinshelwood, proporcional a $\theta_{\ce{CO}}\theta_{\ce{O2}}$, está limitada por la
escasez de oxígeno sobre una superficie abarrotada de CO.
\end{solution}

\begin{solution}{exo:b3:surfaces-catalysis:7}
$\Delta_{\mathrm{ads}}H = -R\ln(8.0/1.0)/(1/300 - 1/350) = -8.314 \times 2.079/\num{4.76e-4} = \qty{-36}{kJ/mol}$.
\end{solution}

\begin{solution}{exo:b3:surfaces-catalysis:8}
$n_{\mathrm m} = 120/22\,414 = \qty{5.35e-3}{mol}$; área $= \num{5.35e-3} \times \num{6.022e23} \times \num{0.162e-18} =
\qty{522}{m^2.g^{-1}}$.
\end{solution}

\begin{solution}{exo:b3:surfaces-catalysis:9}
Langmuir--Hinshelwood: ambos reactivos deben estar adsorbidos en sitios vecinos;
el CO se enlaza fuertemente y, a alta presión, no deja sitio al oxígeno.
\end{solution}

\begin{solution}{exo:b3:surfaces-catalysis:10}
Bajo \omterm{def:b3:surfaces-catalysis:coverage}{grado de recubrimiento}: $100 - 60 = \qty{40}{kJ/mol}$; alto \omterm{def:b3:surfaces-catalysis:coverage}{grado de recubrimiento}: \qty{100}{kJ/mol}. La gráfica de $\ln r$
en función de $1/T$ es pronunciada (pendiente $-100/R$) a baja temperatura, donde la superficie está
llena, y más suave (pendiente $-40/R$) a alta temperatura, donde se vacía: una curva
que se dobla entre ambas.
\end{solution}

\begin{solution}{exo:b3:surfaces-catalysis:11}
Están bloqueados $4 \times 0.10 = 0.40$ de los sitios: queda un 60\,\% de la actividad para un
sitio, y alrededor de $0.60^2 = 36\,\%$ para dos sitios libres adyacentes.
\end{solution}

\begin{solution}{exo:b3:surfaces-catalysis:12}
El intermedio (\qty{-250}{kJ/mol}): el primero enlaza el nitrógeno demasiado débilmente para romper el \ce{N2}
con facilidad, y el tercero demasiado fuertemente, de modo que los átomos de nitrógeno se quedan en la superficie
y la bloquean.
\end{solution}

\begin{solution}{pb:b3:surfaces-catalysis:1}
\textbf{1.} El nitrógeno condensa hacia \qty{77}{K} a presión atmosférica: se alcanza todo el intervalo de
presiones relativas y se forman multicapas.
\textbf{2.} Por debajo de 0.05, la heterogeneidad de la superficie, y por encima de 0.30, la condensación en los
poros, contradicen las hipótesis BET.
\textbf{3.} \num{1.278e-3}, \num{2.410e-3}, \num{3.453e-3}, \num{4.621e-3}, \num{5.659e-3}, \num{6.813e-3} \unit{g.cm^{-3}}.
\textbf{4.} $v_{\mathrm m} = 1/(0.02205 + 0.000180) = \qty{45.0}{cm^3.g^{-1}}$.
\textbf{5.} $c = 1 + 0.02205/0.000180 = 124$: la primera capa se enlaza mucho más fuertemente que
el líquido.
\textbf{6.} Es diminuta comparada con la pendiente multiplicada por $x$, de modo que su error relativo es grande;
pero $v_{\mathrm m}$ depende de la suma de la pendiente y la ordenada, dominada por la pendiente.
\textbf{7.} $45.0/22\,414 = \qty{2.01e-3}{mol.g^{-1}}$.
\textbf{8.} $\num{2.01e-3} \times \num{6.022e23} \times \num{0.162e-18} = \qty{196}{m^2.g^{-1}}$.
\textbf{9.} El soporte de alúmina: el platino expuesto (Parte III) tiene alrededor de \qty{0.9}{m^2.g^{-1}}.
\textbf{10.} Los poros de anchura molecular se llenan por completo a presión muy baja, no capa
a capa, de modo que no puede identificarse ninguna \omterm{def:b3:surfaces-catalysis:coverage}{monocapa}.
\textbf{11.} $2 \times 0.205/22\,414 = \qty{1.83e-5}{mol.g^{-1}}$.
\textbf{12.} $0.0100/195.08 = \qty{5.13e-5}{mol.g^{-1}}$.
\textbf{13.} $D = 1.83/5.13 = 0.357$.
\textbf{14.} $a^3/4 = (0.3923)^3/4 = \qty{0.01509}{nm^3}$.
\textbf{15.} $\frac{\sqrt3}{4}(0.3923)^2 = \qty{0.0666}{nm^2}$, la cara más densa; (100) y (110) dan 0.0770 y
\qty{0.1088}{nm^2}, y la media de las tres densidades de sitios da \qty{0.0807}{nm^2}.
\textbf{16.} Una esfera contiene $\pi d^3/6v_{\mathrm{at}}$ átomos, de los cuales $\pi d^2/a_{\mathrm{at}}$ están en la superficie:
$D = 6v_{\mathrm{at}}/(a_{\mathrm{at}}d)$.
\textbf{17.} $d = 6 \times 0.01509/(0.0807 \times 0.357) = \qty{3.1}{nm}$.
\textbf{18.} Un H por Pt superficial; partículas esféricas de un solo tamaño (el resultado es una
media ponderada por la superficie); todo el platino reducido y accesible; ningún desbordamiento
del hidrógeno hacia el soporte.
\textbf{19.} $\num{2.0e-5}/\num{1.83e-5} = \qty{1.1}{s^{-1}}$.
\textbf{20.} $\num{2.0e-5}/0.0100 = \qty{2.0e-3}{mol.g^{-1}.s^{-1}}$ por gramo de platino.
\textbf{21.} $D = 1.12/10 = 0.112$ en lugar de 0.357: un factor de 3.2.
\textbf{22.} Que la reacción es sensible a la estructura: sus \omterm{def:b3:complex-kinetics:enzyme}{sitios activos} (aristas,
vértices, caras particulares) no son una fracción constante de la superficie.
\textbf{23.} Cuenta los sucesos por \omterm{def:b3:complex-kinetics:enzyme}{sitio activo}, eliminando el efecto de la cantidad de
metal expuesta.
\textbf{24.} \textbf{Diámetro medio de las partículas de platino $\approx \qty{3.1}{nm}$.}
\end{solution}
